\scanpage{71}
$X$ projective lisse \inserted{géom. connexe} sur $k$ corps fini à $q$ éléments. % 1

Par passage à la limite sur les \deleted{quotients} \inserted{coefficients} finis, on % 2
trouve % 3
\[
0\to H^{i-1}(\bar X,\mathbb{Q}_{\ell}(j))_{\pi}
\longrightarrow H^i(\deleted{\unsure{\overline{\phantom X}}}X,\mathbb{Q}_{\ell}(j))
\longrightarrow H^i(\bar X,\mathbb{Q}_{\ell}(j))_{\deleted{\unsure{\pi}}}^{\pi}\to 0 % 4
\]
d'où le résultat % 5
\[
\begin{aligned}
\text{si } i\ne2j,2j+1,\qquad & H^i(X,\mathbb{Q}_{\ell}(j))=0\\ % 6
\text{Si } i=2j\qquad & H^{\deleted{\unsure{i}}2j}(X,\mathbb{Q}_{\ell}(j))
 = H^{2j}(\bar X,\mathbb{Q}_{\ell}(j))^{\pi}\\ % 7
&\hspace{38mm}\rotatebox{90}{$\simeq$}\ \text{Conj. de Tate}\\ % 8
&\hspace{34mm} A^j(X/k)\otimes_{\deleted{\mathbb{Q}}\mathbb{Z}}\mathbb{Q}_{\ell}\\ % 9
\text{si } i=2j+1\qquad & H^{2j+1}(X,\mathbb{Q}_{\ell}(j))
 =H^{2j\deleted{\unsure{+1}}}(\bar X,\mathbb{Q}_{\ell}(j))_{\pi}\\ % 10
&\hspace{38mm}\rotatebox{90}{$\simeq$}\ \text{Tate + action semi-simple}\\ % 11
&\hspace{34mm} A^j(X/k)\otimes_{\deleted{\mathbb{Q}}\mathbb{Z}}\mathbb{Q}_{\ell}
\end{aligned} % 12
\]
Soit maintenant $k_n$ l'extension de degré % 13
$n$ de $k$ dans $\bar k$, et étudions % 14
\[
\varinjlim_n H^i(X_{k_n},\mathbb{Q}_\ell(j)) % 15
\]
On trouve $0$ pour $i\ne2j,2j+1$ ; pour % 16
$i=2j$, on trouve $H^{2j}(\bar X,\mathbb{Q}_{\ell}(j))^{\unsure{oj|of|0j|\mathfrak{g}}}=A^j(\bar X)\otimes_{\deleted{\mathbb{Q}}\mathbb{Z}}\mathbb{Q}_{\ell}$, % 17
enfin pour $i=2j+1$, on trouve également % 18
$A^j(\bar X)\otimes_{\mathbb{Z}}\mathbb{Q}_{\ell}$. \deleted{\unsure{Donc|De}} \unsure{on|$=$}, on trouve \inserted{du pt de vue motifs} % 19
\[
\begin{aligned}
\varinjlim H^i(X_{k_{\unsure{n|\infty}}},\mathbb{Q}(j))
&=\left\{\begin{aligned}
&0 &&\text{si }i\ne2j,2j+1\\ % 20
&A^j(\bar X)\otimes_{\deleted{\mathbb{Q}}\mathbb{Z}}\mathbb{Q}&&\text{si }i=2j\deleted{,2j+1}\\ % 21
&\deleted{\unsure{A^j(\bar X)\otimes\mathbb{Q}}}\ \text{id}&&\text{si }i=2j+1,
\end{aligned}\right.\\ % 22
\rotatebox{90}{$=$}\qquad&\hspace{17mm}\text{mais pas \uline{canoniquement} en}\\ % 23
H^i(\bar X,\mathbb{Q}(j))\qquad&\hspace{17mm}\text{\unsure{donnant|donnons|termes} un }\unsure{\bar x|x|\bar X},\ \bar\pi\,!
\end{aligned} % 24
\]
\margin{\uline{N.B.} \unsure{Remplaçons|Remplacer} $k$ par $k_n$ revient à multiplier l'isom obtenu par $\unsure{n'/n|n/n|1/n}$ ---} % 25
% archive sheet number: 70
